% Diffraction par des objets
Un laser émet une lumière rouge de longueur d'onde $\lambda=\qty{633}{\nm}$. Le
faisceau produit, très directif, a une section pratiquement constante, de
$\qty{1}{\mm}$ de diamètre, et vient frapper un écran en un point $O$.

On interpose sur le trajet du faisceau des plaques percées de trous ou de fentes
de différentes dimensions. On observe l'écran
(Fig.~\ref{fig:montage_diffraction}).

Décrire l'aspect de l'écran lorsque les objets de la Fig.~\ref{fig:ouvertures}
sont interposés.

\begin{minipage}{0.49\linewidth}
    \begin{figure}[H]
        \centering
        \includegraphics[width=\textwidth]{eed31a35-644c-4186-8038-51efc5fb6ca2-3}
        \caption{}
        \label{fig:montage_diffraction}
    \end{figure}
\end{minipage}
\hfill
\begin{minipage}{0.49\linewidth}
\begin{figure}[H]
    \centering
    \includegraphics[width=\textwidth]{eed31a35-644c-4186-8038-51efc5fb6ca2-3(1)}
    \caption{}
    \label{fig:ouvertures}
\end{figure}
\end{minipage}
